\section{\texorpdfstring{$\R^n$}{R n} 的局部线性化与体积几何}\label{sec:02-04}

\subsection{\Frechet{} 微分、链式法则与逆映射}\label{subsec:2-4-1}
\index{Frechet@\Frechet{} 微分}\index{逆函数定理}\index{隐函数定理}

一元函数在一点的导数是一个数；在多元问题中，候选者变成了线性映射。真正需要小心的不是把
偏导数排成矩阵，而是矩阵是否同时控制所有趋近方向。下面这个函数把三种常被混在一起的说法
彻底分开。

\begin{example}[方向导数齐全，函数却不连续]\label{ex:2-4-1-1}
定义
\[
 f(0,0)=0,\qquad
 f(x,y)=\frac{x^3y}{x^6+y^2}\quad ((x,y)\ne(0,0)).
\]
固定 $(a,b)\in\R^2$。若 $b\ne0$，则
\[
 f(ta,tb)=\frac{t^2a^3b}{t^4a^6+b^2},\qquad
 \frac{f(ta,tb)-f(0,0)}{t}\longrightarrow0;
\]
若 $b=0$，上述商恒为零。因此 $f$ 在原点沿每条直线的方向导数都存在且等于零，特别地两个
偏导数均为零。可是沿 $y=x^3$ 有
\[
 f(x,x^3)=\frac{x^6}{2x^6}=\frac12\qquad(x\ne0),
\]
所以 $f$ 在原点不连续。逐方向的极限允许误差随方向改变；一个统一的一阶近似不能允许这种逃逸。
\end{example}

三个具体问题说明统一线性近似所增加的内容。连续偏导能否把坐标方向的资料
合并成同一个线性映射？极坐标映射
$\Phi(r,\theta)=(r\cos\theta,r\sin\theta)$ 在何处有局部逆，而周期性为何仍阻止全局反演？
圆周方程 $x^2+y^2-1=0$ 在何处能解成 $y=g(x)$，又为何在左右端点必须交换坐标？
链式法则、逆函数定理与隐函数定理将依次回答这些问题。

\begin{definition}[\Frechet{} 微分]\label{def:2-4-1-1}
设 $U\subset\R^n$ 为开集，$f:U\to\R^m$，且 $x\in U$。若存在线性映射
$A:\R^n\to\R^m$ 使
\[
 \lim_{\substack{h\to0\\x+h\in U}}
 \frac{|f(x+h)-f(x)-Ah|}{|h|}=0,
 \tag{2.4.1}\label{eq:frechet-definition}
\]
则称 $f$ 在 $x$ \emph{\Frechet{} 可微}，并写 $Df(x)=A$。在标准基下，$Df(x)$ 的矩阵称为
Jacobian 矩阵。
\end{definition}

\begin{lemma}[微分的唯一性与连续性]\label{lemma:2-4-1-derivative-unique-continuous}
满足 \eqref{eq:frechet-definition} 的线性映射至多一个；而且，函数在一点 \Frechet{} 可微便在该点
连续。
\end{lemma}

\begin{proof}
若 $A,B$ 都满足定义，则对每个单位向量 $v$ 以及充分小的非零 $t$，点 $x+tv$ 属于 $U$，并且
\[
 |(A-B)v|
 \le \frac{|f(x+tv)-f(x)-Atv|}{|t|}
    +\frac{|f(x+tv)-f(x)-Btv|}{|t|}.
\]
令 $t\to0$ 得 $(A-B)v=0$。任意非零向量都是单位向量的数乘，故 $A=B$。

再将余项记作 $r(h)=f(x+h)-f(x)-Df(x)h$。由
\[
 |f(x+h)-f(x)|\le \|Df(x)\|\,|h|+|r(h)|
\]
及 $|r(h)|/|h|\to0$，右端随 $h\to0$ 趋于零。
\end{proof}

\begin{definition}[$C^1$ 映射]\label{def:2-4-1-2}
若 $f$ 在开集 $U$ 的每一点 \Frechet{} 可微，且 $x\mapsto Df(x)$ 关于算子范数连续，则写
$f\in C^1(U,\R^m)$。
\end{definition}

\begin{definition}[局部微分同胚]\label{def:2-4-1-3}
若存在 $x$ 的开邻域 $V$ 与 $f(x)$ 的开邻域 $W$，使得 $f:V\to W$ 双射且 $f$、$f^{-1}$
均为 $C^1$，则称 $f$ 在 $x$ 附近给出一个局部 $C^1$ 微分同胚。
\end{definition}

定义中的余项量词正好足以在复合时传递。

\begin{theorem}[链式法则]\label{theorem:2-4-1-1}
设 $U\subset\R^n$ 与 $V\subset\R^m$ 为开集，$f:U\to V$ 在 $x\in U$ 可微，且
$g:V\to\R^p$ 在 $f(x)$ 可微。则 $g\circ f$ 在 $x$ 可微，且
\[
 D(g\circ f)(x)=Dg(f(x))Df(x).
\]
\end{theorem}

\begin{proof}
记 $A=Df(x)$、$B=Dg(f(x))$，并写
\[
 f(x+h)-f(x)=Ah+r(h),\qquad \frac{|r(h)|}{|h|}\to0.
\]
于是存在 $h=0$ 的一个邻域，使
$|f(x+h)-f(x)|\le(\|A\|+1)|h|$。令 $k=f(x+h)-f(x)$。在 $g$ 的可微性公式中写
\[
 g(f(x)+k)-g(f(x))=Bk+s(k),\qquad \frac{|s(k)|}{|k|}\to0.
\]
当 $k=0$ 时取 $s(k)=0$。合并两式得到
\[
 g(f(x+h))-g(f(x))-BAh=Br(h)+s(k).
\]
第一项除以 $|h|$ 后趋于零；第二项满足
如下无除零问题的估计。给定 $\varepsilon>0$，当 $k$ 足够小时（包括 $k=0$，此时 $s(0)=0$）有
$|s(k)|\le\varepsilon|k|$；而 $|k|\le(\|A\|+1)|h|$。所以
\[
 \frac{|s(k)|}{|h|}\le\varepsilon(\|A\|+1).
\]
先令 $h\to0$，再令 $\varepsilon\downarrow0$，得到该商趋于零。这便是以 $BA$ 为微分的
\eqref{eq:frechet-definition}。
\end{proof}

\begin{proposition}[$C^1$ 判据与均值估计]\label{proposition:2-4-1-2}
设 $U\subset\R^n$ 开，$f:U\to\R^m$ 的全部一阶偏导数存在并连续，则 $f\in C^1(U,\R^m)$，
且 $Df(x)$ 是由偏导数组成的 Jacobian 矩阵。若线段 $[x,y]\subset U$，则
\[
 |f(y)-f(x)|
 \le \sup_{z\in[x,y]}\|Df(z)\|\,|y-x|.
 \tag{2.4.2}\label{eq:vector-mean-estimate}
\]
\end{proposition}

\begin{proof}
固定 $x\in U$。取一个以 $x$ 为中心、闭包包含于 $U$ 的开盒。对足够小的
$h=(h_1,\dots,h_n)$，置
\[
 x^{(j)}=x+(h_1,\dots,h_j,0,\dots,0),\qquad x^{(0)}=x.
\]
对 $f$ 的第 $\ell$ 个分量在第 $j$ 段应用一元均值定理，得到位于该坐标线段上的点
$\xi_{j\ell}$，使
\[
 f_\ell(x^{(j)})-f_\ell(x^{(j-1)})
 =\partial_j f_\ell(\xi_{j\ell})h_j.
\]
所有 $\xi_{j\ell}$ 都随 $h\to0$ 趋于 $x$。将坐标段相加并减去
$\sum_j\partial_jf_\ell(x)h_j$。对充分小的 $r>0$ 定义
\[
 \omega_x(r)=
 \max_{\substack{1\le j\le n\\1\le\ell\le m}}
 \sup_{|z-x|\le r}
 |\partial_jf_\ell(z)-\partial_jf_\ell(x)|.
\]
闭球取在前述盒内；有限多个偏导在 $x$ 连续，所以 $\omega_x(r)\to0$。于是
\[
 \left|f_\ell(x+h)-f_\ell(x)-\sum_{j=1}^n\partial_jf_\ell(x)h_j\right|
 \le \omega_x(|h|)\sum_{j=1}^n|h_j|,
\]
由 $\sum_j|h_j|\le\sqrt n\,|h|$，再对 $m$ 个分量使用 Euclidean 范数，余项至多为
$\sqrt{mn}\,\omega_x(|h|)|h|=o(|h|)$，从而得到
\Frechet{} 余项。偏导连续又说明所得矩阵随 $x$ 连续。

现在设 $[x,y]\subset U$。若 $f(y)=f(x)$，估计成立。否则令
\[
 u=\frac{f(y)-f(x)}{|f(y)-f(x)|},\qquad
 \varphi(t)=\langle f(x+t(y-x)),u\rangle .
\]
标量函数 $\varphi$ 在 $[0,1]$ 上连续、在 $(0,1)$ 上可微。一元均值定理给出
$\theta\in(0,1)$，使
\[
 |f(y)-f(x)|=\varphi(1)-\varphi(0)
 =\langle Df(x+\theta(y-x))(y-x),u\rangle.
\]
Cauchy--Schwarz 不等式随即给出 \eqref{eq:vector-mean-estimate}。
\end{proof}

逆函数定理的存在部分归结为一个完备性论证。把它表述在一般完备度量空间中，可以看清欧氏闭球
只是应用场景，而不是原理本身。

\begin{lemma}[压缩映射原理]\label{lemma:2-4-1-2a}
设 $(X,d)$ 是非空完备度量空间，$T:X\to X$ 满足
\[
 d(Tx,Ty)\le qd(x,y)\qquad(x,y\in X)
\]
其中 $0\le q<1$。则 $T$ 有唯一不动点。特别地，完备度量空间中的非空闭子集，或 Banach
空间中被 $T$ 映入自身的非空闭球，都可以作为这里的 $X$。
\end{lemma}

\begin{proof}
任取 $x_0\in X$ 并递归定义 $x_{k+1}=Tx_k$。归纳得到
$d(x_{k+1},x_k)\le q^kd(x_1,x_0)$。当 $p\ge1$ 时，三角不等式给出
\[
 d(x_{k+p},x_k)
 \le\sum_{j=k}^{k+p-1}d(x_{j+1},x_j)
 \le\frac{q^k}{1-q}d(x_1,x_0).
\]
因此 $(x_k)$ 是 Cauchy 列。完备性给出 $x_k\to x_*\in X$。压缩估计还说明 $T$ 连续，故
$Tx_*=\lim_kTx_k=\lim_kx_{k+1}=x_*$。若 $y_*$ 也是不动点，则
$d(x_*,y_*)\le qd(x_*,y_*)$，从 $q<1$ 得 $x_*=y_*$。

若 $F$ 是完备空间的闭子集，则 $F$ 完备；闭球是闭子集。这证明最后一项说明。
\end{proof}

\begin{example}[可算的非线性反演]\label{ex:2-4-1-4}
考虑
\[
 x+\varepsilon\sin x=y,
 \]
并假设 $|\varepsilon|<1$。对固定的 $y$，令
$T_y(x)=y-\varepsilon\sin x$。区间
$I_y=[y-|\varepsilon|,y+|\varepsilon|]$ 满足
\[
 |T_y(x)-y|\le|\varepsilon|,
\]
所以 $T_y(I_y)\subset I_y$。又由一元均值定理，
\[
 |T_y(x)-T_y(x')|\le|\varepsilon|\,|x-x'|.
\]
任取 $x_0\in I_y$，再迭代 $x_{k+1}=y-\varepsilon\sin x_k$；不变性保证所有 $x_k\in I_y$。
压缩映射原理给出唯一解，并给出误差界
\[
 |x_k-x_*|\le\frac{|\varepsilon|^k}{1-|\varepsilon|}|x_1-x_0|.
\]
闭区间的不变性与小于一的压缩常数缺一不可；前者保证迭代始终有定义，后者保证迭代收敛。
\end{example}

\begin{theorem}[逆函数定理]\label{theorem:2-4-1-3}
设 $U\subset\R^n$ 开，$f\in C^1(U,\R^n)$，且 $Df(x_0)$ 可逆。则存在开邻域
$V\ni x_0$ 与 $W\ni f(x_0)$，使 $f:V\to W$ 为 $C^1$ 微分同胚。对 $x\in V$，
\[
 D(f^{-1})(f(x))=Df(x)^{-1}.
 \tag{2.4.3}\label{eq:inverse-derivative}
\]
\end{theorem}

\begin{proof}
令 $A=Df(x_0)$。经过定义域平移和目标空间中的可逆线性变换，问题化为
\[
 x_0=0,\qquad f(0)=0,\qquad Df(0)=I.
 \tag{2.4.4}\label{eq:inverse-normalization}
\]
确切地说，可对 $u\mapsto A^{-1}(f(x_0+u)-f(x_0))$ 证明结论，再把坐标变换复原。

固定 $q\in(0,1)$。由 $Df$ 在 $0$ 连续，可取 $r>0$，使
$\overline B(0,r)\subset U$ 且
\[
 \|Df(u)-I\|\le q\qquad(|u|\le r).
 \tag{2.4.5}\label{eq:inverse-q-bound}
\]
置 $R(u)=f(u)-u$。闭球是凸集，命题 \ref{proposition:2-4-1-2} 应用于 $R$ 给出
\[
 |R(u)-R(v)|\le q|u-v|\qquad(u,v\in\overline B(0,r)).
 \tag{2.4.6}\label{eq:inverse-R-contraction}
\]

令 $W_0=B(0,(1-q)r)$。对每个 $y\in W_0$，在闭球上定义
\[
 T_y(u)=y-R(u).
\]
由于 $R(0)=0$，对 $|u|\le r$ 有
\[
 |T_y(u)|\le |y|+|R(u)-R(0)|<(1-q)r+qr=r.
\]
所以 $T_y$ 把闭球映入自身；\eqref{eq:inverse-R-contraction} 又表明它是 $q$-压缩。存在唯一
$u\in\overline B(0,r)$ 满足 $T_y(u)=u$，而该等式正是 $f(u)=y$。上面的严格不等式还给出
$u\in B(0,r)$。

另一方面，对 $u,v\in\overline B(0,r)$，
\[
 |f(u)-f(v)|
 =|(u-v)+(R(u)-R(v))|
 \ge(1-q)|u-v|.
 \tag{2.4.7}\label{eq:inverse-lower-lipschitz}
\]
故 $f$ 在闭球上单射。令
\[
 V_0=f^{-1}(W_0)\cap B(0,r).
\]
$V_0$ 是包含 $0$ 的开集，以上存在性和唯一性说明 $f:V_0\to W_0$ 双射，而且
\[
 |f^{-1}(y')-f^{-1}(y)|\le(1-q)^{-1}|y'-y|.
 \tag{2.4.8}\label{eq:inverse-lipschitz}
\]
因此逆映射连续。

还需证明逆映射可微，不能预先对它使用链式法则。取 $y=f(u)\in W_0$，令
$y'=f(u')\to y$。由 \eqref{eq:inverse-lipschitz}，$u'-u=O(|y'-y|)$。在 $u$ 处的可微性给出
\[
 y'-y=Df(u)(u'-u)+\rho(u'-u),
 \qquad \frac{|\rho(h)|}{|h|}\to0.
 \tag{2.4.9}\label{eq:inverse-remainder}
\]
由 \eqref{eq:inverse-q-bound}，
$|Df(u)v|\ge(1-q)|v|$；在有限维空间中这说明 $Df(u)$ 可逆，且
$\|Df(u)^{-1}\|\le(1-q)^{-1}$。因而
\[
 \frac{|u'-u-Df(u)^{-1}(y'-y)|}{|y'-y|}
 \le \frac{1}{1-q}\frac{|\rho(u'-u)|}{|u'-u|}
       \frac{|u'-u|}{|y'-y|}\longrightarrow0.
\]
这证明 \eqref{eq:inverse-derivative}。最后，矩阵求逆在可逆矩阵集上连续：若 $B$ 接近可逆矩阵
$C$，置 $\theta=\|C^{-1}(B-C)\|<1$。压缩映射原理应用于
$v\mapsto C^{-1}w-C^{-1}(B-C)v$ 给出 $B^{-1}w$，并给出局部一致的逆范数界；再由
\[
 \|B^{-1}w\|
 \le\|C^{-1}\|\,\|w\|+\theta\|B^{-1}w\|,
 \qquad
 \|B^{-1}\|\le\frac{\|C^{-1}\|}{1-\theta}.
\]
于是 $B\to C$ 时这些逆范数局部一致有界，再由
\[
 B^{-1}-C^{-1}=B^{-1}(C-B)C^{-1}
\]
得 $\|B^{-1}-C^{-1}\|\to0$。因此
$Df(u)^{-1}$ 随 $u$ 连续，\eqref{eq:inverse-lipschitz} 又给 $u=f^{-1}(y)$ 的连续性，故
$f^{-1}\in C^1(W_0)$。

复原 \eqref{eq:inverse-normalization} 前的两次坐标变换，取
$V=x_0+V_0$ 与 $W=f(x_0)+AW_0$，结论成立。
\end{proof}

\begin{example}[极坐标映射]\label{ex:2-4-1-2}
令
\[
 \Phi(r,\theta)=(r\cos\theta,r\sin\theta).
\]
其导数和行列式为
\[
 D\Phi(r,\theta)=
 \begin{pmatrix}
 \cos\theta&-r\sin\theta\\
 \sin\theta&r\cos\theta
 \end{pmatrix},
 \qquad \det D\Phi(r,\theta)=r.
\]
所以在每个 $r_0>0$ 的点附近，逆函数定理给出局部 $C^1$ 逆。$r=0$ 时行列式消失，所有角度
都被送到原点；即使 $r>0$，若角度不限制在长度小于 $2\pi$ 的区间中，$2\pi$ 周期性也破坏全局
单射。局部结论与全局结论在这里可以直接从公式中区分。
\end{example}

\begin{corollary}[隐函数定理]\label{corollary:2-4-1-4}
设 $F:U\subset\R^p\times\R^q\to\R^q$ 为 $C^1$，
$F(x_0,y_0)=0$，且 $D_yF(x_0,y_0)$ 可逆。则存在 $x_0$ 的开邻域 $X$、$y_0$ 的开邻域
$Y$ 以及唯一的 $C^1$ 映射 $g:X\to Y$，使得
\[
 \{(x,y)\in X\times Y:F(x,y)=0\}
 =\{(x,g(x)):x\in X\}.
\]
并且
\[
 Dg(x)=-[D_yF(x,g(x))]^{-1}D_xF(x,g(x)).
 \tag{2.4.10}\label{eq:implicit-derivative}
\]
\end{corollary}

\begin{proof}
定义
\[
 H(x,y)=(x,F(x,y)).
\]
在 $(x_0,y_0)$ 处，
\[
 DH(x_0,y_0)=
 \begin{pmatrix}
 I_p&0\\
 D_xF(x_0,y_0)&D_yF(x_0,y_0)
 \end{pmatrix}.
\]
给定 $(a,b)\in\R^p\times\R^q$，方程
$DH(x_0,y_0)(u,v)=(a,b)$ 有唯一解
\[
 u=a,\qquad
 v=D_yF(x_0,y_0)^{-1}(b-D_xF(x_0,y_0)a),
\]
故该块矩阵可逆。逆函数定理给出 $H$ 的局部 $C^1$ 逆。因为 $H$ 的第一坐标就是 $x$，逆映射
的邻域可显式缩小如下。设定理给出
$H:V_0\to W_0$ 的微分同胚，且 $(x_0,y_0)\in V_0$、$(x_0,0)\in W_0$。先取开邻域
$X_0,Y$ 使 $X_0\times Y\subset V_0$。由 $H^{-1}$ 在 $(x_0,0)$ 连续，再取开邻域
$X_1\ni x_0$ 与 $Z\ni0$，使
\[
 X_1\times Z\subset W_0,
 \qquad H^{-1}(X_1\times Z)\subset X_0\times Y.
\]
令 $X=X_0\cap X_1$。由 $H$ 的第一坐标恒为输入的第一坐标，逆映射在 $X\times Z$ 上必有形式
\[
 H^{-1}(x,z)=(x,h(x,z)).
\]
令 $g(x)=h(x,0)$，便有 $F(x,g(x))=0$。局部逆的唯一性同时说明，缩小邻域后所有零点都以此
形式出现：若 $(x,y)\in X\times Y$ 且 $F(x,y)=0$，则
\[
 H(x,y)=(x,0)=H(x,g(x));
\]
而两点都在 $V_0$，$H|_{V_0}$ 单射，故 $y=g(x)$。

对恒等式 $F(x,g(x))=0$ 使用链式法则，得到
\[
 D_xF(x,g(x))+D_yF(x,g(x))Dg(x)=0.
\]
连续性保证在邻域缩小后 $D_yF(x,g(x))$ 仍可逆，左乘其逆即得
\eqref{eq:implicit-derivative}。
\end{proof}

\begin{example}[圆周在局部是一张图]\label{ex:2-4-1-3}
取 $F(x,y)=x^2+y^2-1$。在圆周上一点 $(x_0,y_0)$ 若 $y_0\ne0$，则
$D_yF(x_0,y_0)=2y_0\ne0$。隐函数定理给出 $y=g(x)$，并由
\eqref{eq:implicit-derivative} 得
\[
 g'(x)=-\frac{x}{g(x)}.
\]
在上半圆和下半圆中，这分别与 $g(x)=\sqrt{1-x^2}$ 和
$g(x)=-\sqrt{1-x^2}$ 的直接求导一致。若 $y_0=0$，不能在该点把圆周写成 $y$ 关于 $x$ 的图，
但 $D_xF=2x_0\ne0$，交换两个坐标后可以写成 $x$ 关于 $y$ 的图。
\end{example}

\begin{figure}[htbp]
  \centering
  \begin{tikzpicture}[scale=0.92]
    \begin{scope}
      \draw[book line] (-2,-1.4) rectangle (0.2,1.4);
      \draw[BookInk!45] (-1.55,-1.4)--(-1.55,1.4) (-0.7,-1.4)--(-0.7,1.4)
                         (-2,-0.45)--(0.2,-0.45) (-2,0.55)--(0.2,0.55);
      \fill[BookAccent!12,draw=BookAccent,semithick] (-1.55,-0.45) rectangle (-0.7,0.55);
      \node[book label] at (-0.9,-1.72) {定义域中的小盒};
      \node[book point] at (-1.12,0.05) {};
      \node[book label,above] at (-1.12,0.12) {$x_0$};
    \end{scope}
    \draw[book accent arrow] (0.45,0.65) -- node[above,book label] {$Df(x_0)$} (2.85,0.65);
    \draw[book dashed arrow] (0.45,-0.65) to[bend right=12]
      node[below,book label] {$f$} (5.75,-0.65);
    \begin{scope}[xshift=4.6cm]
      \fill[BookWarm!12,draw=BookWarm,semithick]
        (-1.45,-0.62)--(-0.55,-0.9)--(0.05,0.45)--(-0.86,0.72)--cycle;
      \node[book label] at (-0.7,-1.72) {线性平行多面体};
    \end{scope}
    \begin{scope}[xshift=7.5cm]
      \fill[BookAccent!10,draw=BookAccent,semithick]
        (-1.5,-0.58) .. controls (-1.16,-0.88) and (-0.78,-0.93) .. (-0.48,-0.8)
        .. controls (-0.15,-0.35) and (0.05,0.15) .. (0.02,0.48)
        .. controls (-0.42,0.78) and (-0.93,0.78) .. (-1.18,0.66)
        .. controls (-1.42,0.25) and (-1.55,-0.22) .. cycle;
      \node[book label] at (-0.72,-1.72) {非线性像};
    \end{scope}
  \end{tikzpicture}
  \caption{在足够小的尺度上，$f$ 的像由 $Df(x_0)$ 的平行多面体加一个低阶误差控制。逆函数定理
  还需要线性部分可逆以及邻域上一致的小误差。}
  \label{fig:2-4-1-linearization}
\end{figure}

逆函数定理只作局部断言。例如 $t\mapsto t^2$ 在 $t=1$ 处导数非零，却在整条实线上不是单射；
边界点也不在上述开集版本的论域内。另一方面，\eqref{eq:inverse-q-bound}--
\eqref{eq:inverse-lower-lipschitz} 留下了比“存在一个逆”更有用的定量信息。下一小节先把其中的线性
映射单独拿出来，核算它对体积的作用。

在常微分方程中，解对初值的局部流若足够光滑，局部反演可以解释短时间内怎样由终点恢复初值；隐函数
定理则用来追踪由方程定义的约束面。随机微分方程的流还需要概率估计和更强的正则性。

\begin{exercise}[微分的唯一性]\label{exercise:2-4-1-1}
不引用引理 \ref{lemma:2-4-1-derivative-unique-continuous}，直接从定义证明：若
$f(x+h)=f(x)+Ah+o(|h|)=f(x)+Bh+o(|h|)$，则 $A=B$。
\end{exercise}

\begin{exercise}[偏导仍然不够]\label{exercise:2-4-1-2}
定义 $u(0,0)=0$，其余点令 $u(x,y)=x^2y/(x^4+y^2)$。计算原点的两个偏导，并研究
$u(x,x^2)$，据此判断连续性和 \Frechet{} 可微性。
\end{exercise}

\begin{exercise}[局部开映射]\label{exercise:2-4-1-3}
设 $f\in C^1(U,\R^n)$ 且 $Df(x_0)$ 可逆。用定理 \ref{theorem:2-4-1-3} 证明：存在
$x_0$ 的邻域 $V$，使 $f(O)$ 对每个相对开集 $O\subset V$ 都是 $f(V)$ 中的相对开集。
\end{exercise}

\begin{exercise}[隐函数导数]\label{exercise:2-4-1-4}
从 $F(x,g(x))=0$ 出发，逐项写出复合映射的微分，并推导
\eqref{eq:implicit-derivative}。说明为什么矩阵乘法的次序不能交换。
\end{exercise}

\subsection{行列式、定向与线性体积缩放}\label{subsec:2-4-2}
\index{行列式}\index{定向}\index{线性换元}

线性映射对小盒的作用可以拆成三种实验：交换坐标，沿某个坐标伸缩，以及把一个坐标的倍数加到
另一个坐标上。前两种的体积变化可以从边长读出；第三种把矩形推成平行多面体，边长相乘的公式不再
直接适用。行列式为三类变换提供同一个乘法因子，但测度公式还必须回答两个问题：剪切为什么
严格保持体积？退化映射为什么把任意集合映入零测集？

行列式最初也可被看作组织线性消元与定向的代数量；在多元微积分中，它获得了新的几何角色：局部
线性化的行列式正是 Jacobian 体积因子。下面先在全局线性映射上严格建立这一解释。

\begin{definition}[行列式的公理刻画]\label{def:2-4-2-1}
在 $(\R^n)^n$ 上唯一满足下列三条性质的函数记作 $\det$：它对每个变量线性；交换两个变量使符号
反转；并且 $\det(e_1,\dots,e_n)=1$。若线性算子 $A$ 的列向量为 $a_1,\dots,a_n$，则
$\det A:=\det(a_1,\dots,a_n)$。
\end{definition}

把 $v_j=\sum_i v_{ij}e_i$ 代入多重线性展开。含有两个相同基向量的项因交换后等于自身的相反数，
故为零；其余项恰由排列 $\sigma\in S_n$ 编号。交换次数决定符号，于是三条性质迫使
\[
 \det(v_1,\dots,v_n)
 =\sum_{\sigma\in S_n}\sgn(\sigma)
   v_{\sigma(1),1}\cdots v_{\sigma(n),n}.
 \tag{2.4.11}\label{eq:det-leibniz-derived}
\]
右端反过来满足三条性质，所以存在性和唯一性都包含在这个有限展开中。

\begin{definition}[定向]\label{def:2-4-2-2}
两个有序基称为同向，如果从一个基到另一个基的过渡矩阵行列式为正。这个关系把全部有序基分成
两个\emph{定向类}。若一个 $C^1$ 微分同胚在所讨论的连通开集上满足 $\det Df>0$，称它保定向；
若 $\det Df<0$，称它反定向。普通 Lebesgue 体积不记录这两个符号，只记录绝对值。
\end{definition}

\begin{proposition}[行列式乘法性]\label{proposition:2-4-2-1}
对任意 $n\times n$ 实矩阵 $A,B$，
\[
 \det(AB)=\det A\det B.
\]
\end{proposition}

\begin{proof}
固定 $A$，考虑
\[
 \Lambda_A(v_1,\dots,v_n)=\det(Av_1,\dots,Av_n).
\]
$\Lambda_A$ 是交替 $n$-线性函数。由定义 \ref{def:2-4-2-1} 的唯一性，任意交替
$n$-线性函数等于其在 $(e_1,\dots,e_n)$ 处的值乘以 $\det$，故
\[
 \Lambda_A(v_1,\dots,v_n)=\det(Ae_1,\dots,Ae_n)\det(v_1,\dots,v_n)
 =\det A\det(v_1,\dots,v_n).
\]
取 $v_1,\dots,v_n$ 为 $B$ 的列向量，左端是 $\det(AB)$，右端是 $\det A\det B$。
\end{proof}

\begin{definition}[单位盒的线性体积比]\label{def:2-4-2-3}
令 $Q_0=[0,1)^n$。对可逆线性映射 $A$，先定义
\[
 J_m(A)=m(AQ_0).
\]
下面的集合公式将证明 $J_m(A)=|\det A|$，并说明以任意正体积盒代替 $Q_0$ 都给出同一体积比。
\end{definition}

首先处理 Gaussian 消元中的三类基本映射。剪切部分的细网格估计是证明中唯一不再把盒送成盒的
地方。

\begin{lemma}[基本线性映射的体积]\label{lemma:2-4-2-elementary-volume}
设 $E\subset\R^n$ Lebesgue 可测。
\begin{enumerate}
  \item 坐标置换保持 $m(E)$；
  \item 若 $D_\lambda$ 只把第 $i$ 个坐标乘以 $\lambda\ne0$，则
        $m(D_\lambda E)=|\lambda|m(E)$；
  \item 若 $S_c$ 把第 $i$ 个坐标替换成 $x_i+cx_j$（$i\ne j$），其余坐标不变，则
        $m(S_cE)=m(E)$。
\end{enumerate}
这些映射还把 Lebesgue 零集的任意子集送到 Lebesgue 零集。
\end{lemma}

\begin{proof}
坐标置换把半开盒送到同体积半开盒。对任意集合的盒覆盖逐个施加映射，再对逆映射重复同一
操作，可得外测度相等。它们是 Borel 双射，因而把 Borel 集送到 Borel 集；结合完成化的 Borel
代表，便得到第一项以及关于零集子集的断言。

对 $D_\lambda$，在 Borel 集上定义 $\nu(B)=m(D_\lambda B)$。由于 $D_\lambda$ 是 Borel 双射，
$\nu$ 是测度。坐标超平面与 $[-M,M]^n$ 的交可放进一个其余边长为 $2M$、法向厚度任意小的盒
中；先令厚度趋于零，再令 $M\in\N$ 取可数并，得到坐标超平面零测。因此当
$\lambda<0$ 时，端点开闭方向的改变不影响测度。于是对任意半开盒 $Q$，
\[
 \nu(Q)=|\lambda|m(Q).
\]
$\nu$ 与 $|\lambda|m$ 都在整数盒耗尽上有限，测度唯一性定理
\ref{theorem:2-2-1-2} 说明二者在全部 Borel 集上相同。若 $E$ Lebesgue 可测，取 Borel 集
$B,N$ 使 $E\mathbin\triangle B\subset N$ 且 $m(N)=0$。Borel 公式给
$m(D_\lambda N)=0$，从
\[
 D_\lambda E\mathbin\triangle D_\lambda B\subset D_\lambda N
\]
得到 $D_\lambda E$ 可测及所需公式。相同论证处理零集的任意子集。

现在考察剪切。设
$Q=\prod_{\ell=1}^n[a_\ell,b_\ell)$，边长为 $d_\ell=b_\ell-a_\ell$。把第 $j$ 个
坐标区间等分成 $N$ 段，长度 $\delta=d_j/N$；若 $Q=\varnothing$，结论平凡，故以下只需考虑
所有 $d_\ell>0$ 的情形。在每个所得薄盒中，$cx_j$ 的振幅不超过
$|c|\delta$。给定 $\varepsilon>0$，该薄盒的像因而包含于一个轴平行盒；后者在第 $i$ 个方向长度为
$d_i+|c|\delta+\varepsilon$，在第 $j$ 个方向长度为 $\delta$，其他方向长度为 $d_\ell$。
这里的任意小外扩处理了半开端点可能不能由同长度半开区间覆盖的问题。于是这 $N$ 个盒给出
\[
 \begin{aligned}
 m^*(S_cQ)
 &\le N(d_i+|c|\delta+\varepsilon)\delta
       \prod_{\ell\ne i,j}d_\ell\\
 &=m(Q)+(|c|\delta+\varepsilon)d_j\prod_{\ell\ne i,j}d_\ell.
 \end{aligned}
\]
先令 $\varepsilon\downarrow0$，再令 $N\to\infty$，得到 $m^*(S_cQ)\le m(Q)$。若任意 $A\subset\R^n$ 被盒列 $(Q_k)$ 覆盖，
则
\[
 m^*(S_cA)\le\sum_km^*(S_cQ_k)\le\sum_km(Q_k).
\]
对所有覆盖取下确界，$m^*(S_cA)\le m^*(A)$。将同一结论用于逆剪切 $S_{-c}$ 和集合 $S_cA$，
得到反向不等式，故
\[
 m^*(S_cA)=m^*(A)\qquad(A\subset\R^n).
 \tag{2.4.12}\label{eq:shear-outer-measure}
\]
若 $E\mathbin\triangle B$ 包含于 Borel 零集 $N$，则 $S_cB$ 是 Borel 集，
\eqref{eq:shear-outer-measure} 说明 $S_cN$ 零测，于是 $S_cE$ 可测且测度不变。零集子集的结论
也直接来自 \eqref{eq:shear-outer-measure}。
\end{proof}

\begin{example}[剪切]\label{ex:2-4-2-1}
在 $\R^2$ 中令 $A(x,y)=(x+cy,y)$。单位方形被送到由四个顶点
\[
 (0,0),\qquad(1,0),\qquad(c,1),\qquad(1+c,1)
\]
张成的平行四边形。上一证明把它按 $N$ 条水平薄带覆盖，总面积
至多 $1+|c|/N$；逆剪切给反向估计，所以其面积等于 $1$。矩阵
\[
 \begin{pmatrix}1&c\\0&1\end{pmatrix}
\]
的行列式也等于 $1$。
\end{example}

\begin{example}[反射]\label{ex:2-4-2-2}
映射
\[
 A(x_1,x_2,\dots,x_n)=(-x_1,x_2,\dots,x_n)
\]
把每个盒的第一坐标区间翻转，体积不变；其矩阵行列式为 $-1$。这项计算说明无符号体积公式若写成
$\det A$，会把非负测度错误地变成负数。
\end{example}

退化矩阵需要另一项覆盖计算。它不能由“把伸缩因子取成零”直接推出，因为尚未知道测度公式对极限
稳定。

\begin{lemma}[低维线性子空间零测]\label{lemma:2-4-2-lower-dimensional-null}
若 $V\subset\R^n$ 是维数 $k<n$ 的线性子空间，则 $m_n(V)=0$。因此 $V$ 的每个子集都是
Lebesgue 可测的零集。
\end{lemma}

\begin{proof}
\(k=0\) 时 \(V=\{0\}\)，而单点可放入边长任意小的 \(n\) 维盒，结论成立。以下设
\(1\le k<n\)。
从坐标投影中选择 $k$ 个坐标，使相应投影 $\pi:V\to\R^k$ 为单射；这是因为 $V$ 的一组基所成
$n\times k$ 矩阵有一个非零的 $k\times k$ 子式。有限维线性代数给出常数 $C\ge1$，使
\[
 |v|\le C|\pi v|\qquad(v\in V).
 \tag{2.4.13}\label{eq:subspace-coordinate-bound}
\]
固定 $M>0$，考察 $V_M=V\cap[-M,M]^n$。集合 $\pi(V_M)$ 位于 $[-M,M]^k$ 中。用至多
$(2M/\varepsilon+2)^k$ 个边长 $\varepsilon$ 的 $k$ 维盒覆盖这个参数盒。由
\eqref{eq:subspace-coordinate-bound}，每个参数小盒对应的 $V_M$ 部分直径至多
$C\sqrt{k}\,\varepsilon$，因而可放进一个边长
$2C\sqrt{k}\,\varepsilon$ 的 $n$ 维坐标盒。于是存在只依赖 $M,n,C$ 的常数 $C_M$，使
\[
\begin{aligned}
 m_n^*(V_M)
 &\le\left(\frac{2M}{\varepsilon}+2\right)^k
       (2C\sqrt{k}\,\varepsilon)^n\\
 &=(2M+2\varepsilon)^k(2C\sqrt{k})^n
       \varepsilon^{n-k}.
\end{aligned}
\]
令 $\varepsilon\downarrow0$ 得 $m_n^*(V_M)=0$。最后
$V=\bigcup_{M\in\N}V_M$，外测度的可数次可加性给 $m_n^*(V)=0$；Lebesgue 测度的完备性处理
所有子集。
\end{proof}

\begin{example}[奇异映射]\label{ex:2-4-2-3}
若 $A:\R^n\to\R^n$ 的秩为 $k<n$，则 $A(\R^n)$ 是一个 $k$ 维线性子空间。对任意
$E\subset\R^n$，
\[
 AE\subset A(\R^n),
\]
所以上一引理给 $m_n^*(AE)=0$。这也包括 $E$ 无界或不可测的情形。
\end{example}

\begin{theorem}[线性换元的集合公式]\label{theorem:2-4-2-2}
设 $A:\R^n\to\R^n$ 线性，$E\subset\R^n$ Lebesgue 可测，则 $AE$ Lebesgue 可测，且
\[
 m(AE)=|\det A|m(E).
 \tag{2.4.14}\label{eq:linear-change-set}
\]
若 $A$ 不可逆，右端约定为零；特别地，当 $m(E)=\infty$ 时也约定
$0\cdot\infty=0$。
\end{theorem}

\begin{proof}
若 $A$ 奇异，则 $\det A=0$，例 \ref{ex:2-4-2-3} 已给出 $AE$ 可测且零测。

设 $A$ 可逆。Gaussian 消元用有限次行交换、非零行伸缩和把一行的倍数加到另一行，把 $A$ 化为
单位矩阵。等价地，$A$ 是有限个坐标置换、坐标伸缩和剪切矩阵的乘积。引理
\ref{lemma:2-4-2-elementary-volume} 可以依次应用于中间像集，得到
\[
 m(AE)=\left(\prod_{r=1}^N c_r\right)m(E),
\]
其中置换、反射和剪切的 $c_r=1$，坐标伸缩的 $c_r$ 是相应伸缩系数的绝对值。另一方面，命题
\ref{proposition:2-4-2-1} 逐项应用，结合三类基本矩阵的直接行列式计算，给出
\[
 |\det A|=\prod_{r=1}^Nc_r.
\]
两式合并即为 \eqref{eq:linear-change-set}。每一步都保持 Lebesgue 可测性，所以中间像与最终像
均可测。
\end{proof}

特别地，对每个正体积半开盒 $Q$，定理给出
\[
 \frac{m(AQ)}{m(Q)}=|\det A|,
\]
从而定义 \ref{def:2-4-2-3} 中的 $J_m(A)$ 的确等于 $|\det A|$，且不依赖所选盒。

\begin{proposition}[奇异值体积公式]\label{proposition:2-4-2-3}
若 $A=U\Sigma V^T$ 是实矩阵的奇异值分解，奇异值为 $s_1,\dots,s_n\ge0$，则
\[
 |\det A|=\prod_{j=1}^ns_j.
\]
单位球在 $A$ 下的像是主半轴长度为 $s_1,\dots,s_n$ 的椭球；若某个 $s_j=0$，这个椭球退化到
低维子空间。
\end{proposition}

\begin{proof}
正交性给 $U^TU=I$ 与 $V^TV=I$。取行列式并使用乘法性，
\[
 (\det U)^2=(\det V)^2=1,
\]
故 $|\det U|=|\det V|=1$。又
$\Sigma=\operatorname{diag}(s_1,\dots,s_n)$，所以
\[
 |\det A|=|\det U|\,|\det\Sigma|\,|\det V^T|
 =\prod_{j=1}^ns_j.
\]
$V^T$ 和 $U$ 都保持欧氏距离，$\Sigma$ 则把第 $j$ 个坐标方向伸缩 $s_j$ 倍；依次作用于单位球便
得到所述椭球。
\end{proof}

\begin{figure}[htbp]
  \centering
  \begin{tikzpicture}[scale=0.88]
    \begin{scope}
      \draw[book line] (-1,-1) rectangle (1,1);
      \draw[BookInk!35] (-1,0)--(1,0) (0,-1)--(0,1);
      \node[book label] at (0,-1.35) {单位盒};
    \end{scope}
    \draw[book accent arrow] (1.25,0) -- node[above,book label] {旋转} (2.35,0);
    \begin{scope}[xshift=3.65cm,rotate=20]
      \draw[book line] (-1,-1) rectangle (1,1);
    \end{scope}
    \node[book label] at (3.65,-1.35) {$|\det|=1$};
    \draw[book accent arrow] (4.9,0) -- node[above,book label] {轴向伸缩} (6.15,0);
    \begin{scope}[xshift=7.45cm]
      \draw[book line] (-1.35,-0.7) rectangle (1.35,0.7);
      \node[book label] at (0,-1.35) {乘 $s_1s_2$};
    \end{scope}
    \draw[book accent arrow] (8.95,0) -- node[above,book label] {剪切} (10.1,0);
    \begin{scope}[xshift=11.45cm]
      \draw[book line] (-1.55,-0.7)--(1.15,-0.7)--(1.55,0.7)--(-1.15,0.7)--cycle;
      \node[book label] at (0,-1.35) {体积不再改变};
    \end{scope}
  \end{tikzpicture}
  \caption{线性体积变化可以分解为正交变换、轴向伸缩与剪切。只有轴向伸缩贡献绝对行列式的大小。}
  \label{fig:2-4-2-linear-volume}
\end{figure}

\begin{proposition}[Jacobian 的局部体积比]\label{corollary:2-4-2-4}
设 $f$ 在 $x$ 的邻域为 $C^1$，且 $Df(x)$ 可逆。逆函数定理保证 $f$ 在更小邻域内单射。对固定的
有界正体积盒 $Q$，其边界零测，则
\[
 \frac{m(f(x+rQ))}{r^nm(Q)}\longrightarrow |\det Df(x)|
 \qquad(r\downarrow0).
 \tag{2.4.15}\label{eq:future-local-jacobian-volume}
\]
这是一个明确延后证明、且本章后续不调用的前瞻命题。证明见\cref{theorem:3-2-2-2}：第~3~章先从
本节已经证明的全局线性体积公式独立建立近恒等映射的局部体积夹逼，再证明集合换元公式；该证明不使用
本命题，因而不存在循环依赖。集合换元公式随后把分子写成 \(J_f\) 在 \(x+rQ\) 上的积分，而
\(J_f\) 的连续性给出极限。若只假设 \(f\) 在 \(x\) 一点可微，则邻域单射性和像集无重叠性都没有
保证。
\end{proposition}

行列式的符号适合记录定向，测度则只能使用其绝对值。非线性映射还可能在不同点把质量送到同一处，
因此线性集合公式并不自动推出非线性换元公式。下一小节只处理一种跨所有尺度的距离控制，并证明它
足以保住零测集。

Gaussian 分布在线性映射下的协方差会按 $A\Sigma A^{T}$ 变换；在有密度的非退化情形中，归一化常数
相应出现 $|\det A|$。Fourier 变换和卷积的线性缩放也会产生同一行列式因子。期望、密度、Fourier
变换与卷积将在积分理论中定义，并以集合公式 \eqref{eq:linear-change-set} 为线性换元的基础。

\begin{exercise}[剪切的外测度版本]\label{exercise:2-4-2-1}
仿照 \eqref{eq:shear-outer-measure} 的证明，写出三维映射
$(x,y,z)\mapsto(x+ay+bz,y,z)$ 的细盒覆盖，并证明它保持每个集合的 Lebesgue 外测度。
\end{exercise}

\begin{exercise}[仿射子空间]\label{exercise:2-4-2-2}
用引理 \ref{lemma:2-4-2-lower-dimensional-null} 证明：$\R^n$ 中每个维数小于 $n$ 的仿射子空间
都是 Lebesgue 零集。
\end{exercise}

\begin{exercise}[椭球体积]\label{exercise:2-4-2-3}
设 $B=\{x:|x|<1\}$，$A$ 可逆。证明
$m(AB)=|\det A|m(B)$，再用奇异值分解把 $AB$ 写成一个主轴椭球并核对其体积因子。
\end{exercise}

\begin{exercise}[有向面积与普通面积]\label{exercise:2-4-2-4}
计算反射 $(x,y)\mapsto(-x,y)$ 对有序基 $(e_1,e_2)$ 的作用。分别说明它怎样改变平行四边形的
有向面积和 Lebesgue 面积。
\end{exercise}

\subsection{Lipschitz 几何、零测集与可求长曲线}\label{subsec:2-4-3}
\index{Lipschitz 映射}\index{Hausdorff 测度}\index{可求长曲线}

连续性只要求很近的点被送到很近的点，却不规定两种“很近”以什么速度对应。下面两个构造展示同一
缺口的两种强度：第一个把零测集映成一个区间，第二个把一条参数区间映满一个正面积方形。它们都
连续；失败发生在尺度估计上。

\begin{example}[Cantor 函数不是 Lipschitz]\label{ex:2-4-3-1}
从 $C_0=[0,1]$ 开始，每一步从每个闭区间删去开中三分之一，记剩余的 $2^k$ 个闭区间之并为
$C_k$，并令 $C=\bigcap_kC_k$。因为
\[
 m(C_k)=2^k3^{-k}=\left(\frac23\right)^k,
\]
而 $C_k\downarrow C$ 且 $m(C_0)=1<\infty$，测度从上连续给
\[
 m(C)=\lim_{k\to\infty}m(C_k)
 =\lim_{k\to\infty}\left(\frac23\right)^k=0.
\]

每个 $x\in C$ 可以选择一个只含数字 $0,2$ 的三进制展开
$x=\sum_{j\ge1}a_j3^{-j}$。定义
\[
 F(x)=\sum_{j\ge1}\frac{a_j}{2}\,2^{-j}\qquad(x\in C),
 \tag{2.4.16}\label{eq:cantor-digit-map}
\]
并在 $[0,1]\setminus C$ 的每个被删开区间上取与两个端点相同的常值。对有两种三进制展开的
端点，选择只含 $0,2$ 的那一种；每个被删区间的两个端点经 \eqref{eq:cantor-digit-map} 恰好得到
同一数值，因此常值延拓无歧义。

先核对单调性。若 $x<y$ 都在 $C$，在它们只含 $0,2$ 的展开首个不同位置 $l$，必有
$a_l(x)=0$、$a_l(y)=2$。因此
\[
 F(y)-F(x)
 \ge2^{-l}-\sum_{j>l}2^{-j}=0.
\]
等号只可能对应某个被删区间的左右端点；在该区间上作常值延拓仍保持次序。因此延拓后的 $F$ 在
$[0,1]$ 上非减。

第 $k$ 层每个基本闭区间上，前 $k$ 个三进制数字固定，故 $F$ 的振幅至多
$\sum_{j>k}2^{-j}=2^{-k}$；被删区间上的振幅为零。

令 $\alpha=\log2/\log3$。若
$3^{-(k+1)}<|x-y|\le3^{-k}$，不同第 $k$ 层基本区间之间的间隙长度至少为 $3^{-k}$，故
$[x,y]$ 至多穿过两个这一级基本区间的有效部分。因此
\[
 |F(x)-F(y)|\le2\cdot2^{-k}\le4|x-y|^\alpha.
\]
这同时严格证明 $F$ 连续，并给出 \Holder{}-$\alpha$ 估计。另一方面，取
$x_k=0$、$y_k=3^{-k}$。数 $3^{-k}$ 有展开
$0.\underbrace{00\cdots0}_{k\text{ 位}}222\cdots$，因而
\[
 |y_k-x_k|=3^{-k},\qquad
 |F(y_k)-F(x_k)|=2^{-k},\qquad
 \frac{|F(y_k)-F(x_k)|}{|y_k-x_k|}=\left(\frac32\right)^k\to\infty.
\]
不存在统一的线性距离常数。最后，任取 $u\in[0,1]$ 的二进制展开
$u=\sum b_j2^{-j}$，把数字 $b_j$ 换成三进制数字 $2b_j$ 便得到 $x\in C$ 且 $F(x)=u$。
于是 $F(C)=[0,1]$：连续映射确实能把零测集送成正测度集。
\end{example}

第二个构造不借用“存在空间填充曲线”这一结论。我们把四个相似变换写出来，让满射从一致极限中
产生。对 $Q=[0,1]^2$ 定义
\[
 \begin{array}{ll}
 S_0(u,v)=(v/2,u/2),
 &S_1(u,v)=(u/2,v/2+1/2),\\[2mm]
 S_2(u,v)=(u/2+1/2,v/2+1/2),
 &S_3(u,v)=(1-v/2,1/2-u/2).
 \end{array}
 \tag{2.4.17}\label{eq:hilbert-similarities}
\]
它们以比例 $1/2$ 把 $Q$ 分别送到左下、左上、右上、右下四个小方形，并安排相邻入口与出口
相接。

\begin{example}[逐层构造的 Peano 曲线]\label{ex:2-4-3-3}
若连续折线 $P:[0,1]\to Q$ 满足 $P(0)=(0,0)$、$P(1)=(1,0)$，定义新折线
\[
 (\mathcal TP)(t)=
 \begin{cases}
 S_0(P(4t)),&0\le t\le1/4,\\
 S_1(P(4t-1)),&1/4\le t\le1/2,\\
 S_2(P(4t-2)),&1/2\le t\le3/4,\\
 S_3(P(4t-3)),&3/4\le t\le1.
 \end{cases}
 \tag{2.4.18}\label{eq:hilbert-operator}
\]
由 \eqref{eq:hilbert-similarities} 直接计算四段的相接值依次为
$(0,0),(0,1/2),(1/2,1/2),(1,1/2),(1,0)$，故 $\mathcal TP$ 连续并保持两个端点。
从 $P_0(t)=(t,0)$ 开始，递归令 $P_{k+1}=\mathcal TP_k$。

每个 $S_i$ 都是相似比 $1/2$ 的映射，所以对任意两条满足端点条件的曲线 $P,R$，
\[
 \|\mathcal TP-\mathcal TR\|_\infty
 \le\frac12\|P-R\|_\infty.
\]
于是
\[
 \|P_{k+1}-P_k\|_\infty
 \le2^{-k}\|P_1-P_0\|_\infty.
\]
对 $m>k$，三角不等式进一步给出
\[
 \|P_m-P_k\|_\infty
 \le\sum_{j=k}^{m-1}\|P_{j+1}-P_j\|_\infty
 \le 2^{1-k}\|P_1-P_0\|_\infty.
\]
故 $(P_k)$ 在一致范数下是 Cauchy 列。$\R^2$ 完备使逐点极限
$P(t)=\lim_kP_k(t)$ 存在；同一尾和与 $k\to\infty$ 给出一致收敛，因此 $P$ 连续。

还需证明极限填满方形。把 $[0,1]$ 按四进制分成 $4^k$ 个闭参数小区间；对每个长度为 $k$ 的词
$w=(i_1,\dots,i_k)\in\{0,1,2,3\}^k$，相应小区间在 $P_k$ 下落入
\[
 Q_w=S_{i_1}\circ\cdots\circ S_{i_k}(Q).
\]
这些 $Q_w$ 恰是 $2^k\times2^k$ dyadic 方格中的全部小方形，每个只出现一次。归纳如下：$k=1$
时这是 \eqref{eq:hilbert-similarities} 的四个像。若长度 $k$ 的词已给出全部父方形，则固定前缀
$w$ 后，四个后继 $wi$ 的像为
\[
 Q_{wi}=S_w(S_i(Q)),\qquad S_w=S_{i_1}\circ\cdots\circ S_{i_k};
\]
它们是 $Q_w$ 的四个 dyadic 子方形。不同前缀属于内部不交的不同父方形，所以不会重复；每个父方形
又被这四个子方形覆盖，归纳完成。四进制参数区间也按同样的前缀规则分成四个嵌套子区间，故词与
参数区间、像方形始终相配。

给定 $z\in Q$，逐层选取包含 $z$ 的嵌套小方形
$Q_{w_k}$，并选取对应的嵌套参数区间 $I_{w_k}$。它们的长度为 $4^{-k}$，故交集只含一点 $t$。
对 $m\ge k$，递归构造保证 $P_m(t)\in Q_{w_k}$。令 $m\to\infty$，再利用
$\diam Q_{w_k}=\sqrt2\,2^{-k}\to0$，得到 $P(t)=z$。因此 $P([0,1])=Q$。

这条曲线不可能满足 $|P(s)-P(t)|\le L|s-t|$。若满足，把 $[0,1]$ 等分成 $N$ 段；每段的像
直径至多 $L/N$。给定 $\eta>0$，每段像可放进边长 $L/N+\eta$ 的半开坐标方形。于是
\[
 1=m(Q)=m^*(P([0,1]))
 \le N\left(\frac LN+\eta\right)^2.
\]
先令 $\eta\downarrow0$ 得 $1\le L^2/N$，再令 $N\to\infty$ 产生矛盾。逐层折线、一致极限与嵌套方格分别保证了可计算性、连续性和满射性。
\end{example}

\begin{definition}[Lipschitz 映射]\label{def:2-4-3-1}
设 $E\subset\R^n$。若 $f:E\to\R^m$ 满足
\[
 |f(x)-f(y)|\le L|x-y|\qquad(x,y\in E),
 \tag{2.4.19}\label{eq:lipschitz-definition}
\]
则称 $f$ 为 $L$-Lipschitz 映射。若定义域的每一点都有一个相对邻域，使 $f$ 在该邻域满足某个
Lipschitz 估计，则称 $f$ 局部 Lipschitz。
\end{definition}

\begin{lemma}[紧集上的统一 Lipschitz 常数]\label{lemma:2-4-3-compact-lipschitz}
若 $f:E\to\R^m$ 局部 Lipschitz，$K\subset E$ 为紧集，则 $f|_K$ 是 Lipschitz 映射。
\end{lemma}

\begin{proof}
若 $K=\varnothing$，结论平凡。以下设 $K\ne\varnothing$。
取有限个相对开集 $U_1,\dots,U_N$ 覆盖 $K$，使 $f$ 在 $U_j$ 上为 $L_j$-Lipschitz。紧度量空间
$K$ 的 Lebesgue 数引理给出 $\delta>0$，使 $K$ 中直径小于 $\delta$ 的集合包含于某个 $U_j$。
因此 $x,y\in K$ 且 $|x-y|<\delta$ 时，
$|f(x)-f(y)|\le L_0|x-y|$，其中 $L_0=\max_jL_j$。局部 Lipschitz 蕴含连续，故
$M=\sup_{x\in K}|f(x)|<\infty$。若 $|x-y|\ge\delta$，则
\[
 |f(x)-f(y)|\le2M\le(2M/\delta)|x-y|.
\]
取 $\max\{L_0,2M/\delta\}$ 即得统一常数。
\end{proof}

线性距离控制还可以跨维数表达。为此只引入本书随后真正会用到的覆盖形式。

\begin{definition}[Hausdorff 内容与测度]\label{def:2-4-3-2}
设 $s\ge0$、$\delta>0$。对度量空间中的 $A$，定义
\[
 \mathcal H^s_\delta(A)=
 \inf\left\{\sum_{i\in I}(\diam U_i)^s:
 A\subset\bigcup_{i\in I}U_i,\ I\text{ 有限或可数},\ \diam U_i\le\delta\right\},
 \tag{2.4.20}\label{eq:hausdorff-delta}
\]
以及
\[
 \mathcal H^s(A)=\lim_{\delta\downarrow0}\mathcal H^s_\delta(A).
\]
删去直径限制得到的下确界记作 $\mathcal H^s_\infty(A)$，称为 Hausdorff 内容。本书采用不含
额外规范化常数的约定。空指标族以成本 $0$ 覆盖空集；$s=0$ 时约定每个非空覆盖块的成本为 $1$、
空覆盖块的成本为 $0$，于是 $\mathcal H^0$ 是计数测度。
\end{definition}

最后一句可由定义直接证明。若 $A$ 含 $N\ge2$ 个不同点 $x_1,\dots,x_N$，令
\[
 d_N=\min_{i\ne j}d(x_i,x_j)>0.
\]
当 $0<\delta<d_N$ 时，每个直径不超过 $\delta$ 的覆盖块至多含一个 $x_i$，所以任一 $A$ 的容许
覆盖至少有 $N$ 个非空块，$\mathcal H^0_\delta(A)\ge N$。若 $A$ 无限，对每个 $N$ 应用此式得到
$\mathcal H^0(A)=\infty$；若 $A$ 有恰好 $N\ge2$ 个点，单点覆盖给反向不等式，故
$\mathcal H^0(A)=N$。单点集的每个覆盖至少含一个非空块且单点覆盖可取，故其值为 $1$；空集由
空覆盖取值 $0$。

\begin{proposition}[Hausdorff 测度是 Borel 测度]\label{proposition:2-4-3-hausdorff-borel}
对每个 $s\ge0$，$\mathcal H^s$ 是度量外测度，因而所有 Borel 集都
$\mathcal H^s$-可测。
\end{proposition}

\begin{proof}
固定 $\delta$，在所有直径不超过 $\delta$ 的集合上取成本 $(\diam U)^s$。覆盖下确界引理
\ref{lemma:2-1-2-1} 说明 $\mathcal H^s_\delta$ 是外测度。随着 $\delta\downarrow0$，这些量单调
增加。对集合列 $(A_j)$，
\[
 \mathcal H^s_\delta\Bigl(\bigcup_jA_j\Bigr)
 \le\sum_j\mathcal H^s_\delta(A_j)
 \le\sum_j\mathcal H^s(A_j).
\]
令 $\delta\downarrow0$，得到 $\mathcal H^s$ 的可数次可加不等式；空集条件与单调性同样从
$\mathcal H^s_\delta$ 传给极限。因此 $\mathcal H^s$ 是外测度。

若 $\dist(A,B)=d>0$，取 $0<\delta<d$。直径不超过 $\delta$ 的覆盖块不可能同时遇到 $A$ 和
$B$。任意一个 $A\cup B$ 的容许覆盖可分成覆盖 $A$ 与覆盖 $B$ 的两个子族，所以
\[
 \mathcal H^s_\delta(A\cup B)
 \ge \mathcal H^s_\delta(A)+\mathcal H^s_\delta(B).
\]
反向不等式来自外测度次可加。令 $\delta\downarrow0$ 得正距离可加性，故
$\mathcal H^s$ 是度量外测度。命题 \ref{proposition:2-1-2-3} 随即给出 Borel
可测性。
\end{proof}

\begin{proposition}[覆盖畸变估计]\label{proposition:2-4-3-1}
若 $f:E\subset\R^n\to\R^m$ 为 $L$-Lipschitz，$A\subset E$，则
\[
 \mathcal H^s(f(A))\le L^s\mathcal H^s(A)\qquad(s>0).
 \tag{2.4.21}\label{eq:hausdorff-lipschitz}
\]
$s=0$ 时，相应结论是像集的点数不超过原集合的点数。若 $m=n$，还有
\[
 m_n^*(f(A))\le n^{n/2}L^n m_n^*(A).
 \tag{2.4.22}\label{eq:lebesgue-lipschitz-distortion}
\]
\end{proposition}

\begin{proof}
对 $U\subset E$，定义直接给出
\[
 \diam f(U)\le L\diam U.
\]
若 $L=0$，像集至多一个点；单点的 $\mathcal H^s$（$s>0$）与 $n$ 维 Lebesgue 外测度都为零，
所以两式均成立。以下设 $L>0$。取 $A$ 的一个直径不超过
$\delta$ 的覆盖；把每个覆盖块与 $E$ 相交后再施加 $f$，得到 $f(A)$ 的直径不超过 $L\delta$ 的
覆盖，且
\[
 \sum_i(\diam f(U_i\cap E))^s
 \le L^s\sum_i(\diam U_i)^s.
\]
先取下确界得到
$\mathcal H^s_{L\delta}(f(A))\le L^s\mathcal H^s_\delta(A)$，再令 $\delta\downarrow0$，得到
\eqref{eq:hausdorff-lipschitz}。$s=0$ 的说法来自函数不能增加集合基数。

证明 \eqref{eq:lebesgue-lipschitz-distortion} 时，若右端无穷则无需再证。否则给定
$\varepsilon>0$，利用命题 \ref{proposition:2-1-3-1a} 取边长为 $r_i$ 的任意细立方体 $Q_i$，使
\[
 A\subset\bigcup_iQ_i,
 \qquad \sum_ir_i^n<m_n^*(A)+\varepsilon.
\]
这里必须考察 $f(Q_i\cap E)$，因为 $f$ 未必定义在整个 $Q_i$ 上。该像集的直径至多
$L\sqrt n\,r_i$。每个坐标函数在该像集上的振幅不超过这个直径；将各坐标值域稍作任意小的
扩大即可得到半开坐标盒。对可数族必须分配可求和误差：再给定 $\eta>0$，对第 $i$ 个非空像集
选盒 $R_i$，使
\[
 |R_i|<(L\sqrt n\,r_i)^n+\eta2^{-i}.
\]
于是
\[
 m_n^*(f(A))
 \le\sum_i|R_i|
 \le n^{n/2}L^n(m_n^*(A)+\varepsilon)+\eta.
\]
先令 $\eta\downarrow0$，再令 $\varepsilon\downarrow0$。
\end{proof}

\begin{theorem}[零测集保持]\label{theorem:2-4-3-2}
若 $f:E\subset\R^n\to\R^n$ 为 Lipschitz 映射，且 $N\subset E$ 满足
$m_n^*(N)=0$，则
\[
 m_n^*(f(N))=0.
\]
特别地，$f(N)$ 自动 Lebesgue 可测且测度为零。
\end{theorem}

\begin{proof}
将 $A=N$ 代入 \eqref{eq:lebesgue-lipschitz-distortion}，得到
\[
 0\le m_n^*(f(N))\le n^{n/2}L^n m_n^*(N)=0.
\]
令 $S=f(N)$。对任意 $A\subset\R^n$，由 $m_n^*(A\cap S)=0$、次可加性和单调性，
\[
 m_n^*(A)
 \le m_n^*(A\setminus S)+m_n^*(A\cap S)
 =m_n^*(A\setminus S)
 \le m_n^*(A).
\]
所以等号成立，$S$ 直接满足 \Caratheodory{} 条件且测度为零。
这个论证没有假设连续像是 Borel 集。
\end{proof}

\begin{example}[Lipschitz 图]\label{ex:2-4-3-2}
设 $n\ge2$，$g:\R^{n-1}\to\R$ 为 $L$-Lipschitz，并记
\[
 \Gamma=\{(x,g(x)):x\in\R^{n-1}\}.
\]
固定 $M$，把 $[-M,M]^{n-1}$ 分成边长不超过 $\delta$ 的小盒；盒数至多
$(2M/\delta+2)^{n-1}$。在每个底盒 $Q$ 中选一点 $x_Q$。若 $x\in Q$，则
\[
 |g(x)-g(x_Q)|\le L\sqrt{n-1}\,\delta.
\]
因此 $Q$ 上方的图可放进一个底边长不超过 $\delta$、高度不超过
$(2L\sqrt{n-1}+1)\delta$ 的 $n$ 维盒。所有这些盒的总体积不超过
\[
 (2M/\delta+2)^{n-1}(2L\sqrt{n-1}+1)\delta^n,
\]
它随 $\delta\downarrow0$ 趋于零，故该图块的 $m_n^*$ 为零。又因 $g$ Lipschitz 从而连续，
\[
 \Gamma=\{(x,t)\in\R^{n-1}\times\R:t-g(x)=0\}
\]
是连续函数 $(x,t)\mapsto t-g(x)$ 的零水平闭集，因而 Borel 可测。故
$m_n(\Gamma\cap([-M,M]^{n-1}\times\R))=0$。再令 $M\in\N$ 并取可数并，得到
$m_n(\Gamma)=0$。

另一方面，参数映射 $x\mapsto(x,g(x))$ 是 $\sqrt{1+L^2}$-Lipschitz；
\eqref{eq:hausdorff-lipschitz} 表明它自然受 $(n-1)$ 维 Hausdorff 测度控制。这不与 $n$ 维零测
矛盾，而是在区分两种维数。
\end{example}

微分提供点态的一阶信息，Lipschitz 条件提供全尺度信息。两者之间的逆向联系由下述经典定理给出。

\begin{remark}[外部背景：Rademacher 定理]\label{theorem:2-4-3-3}
设 $U\subset\R^n$ 开，$f:U\to\R^m$ 局部 Lipschitz，则 $f$ 在 Lebesgue 几乎处处
\Frechet{} 可微。
\end{remark}

Rademacher 定理的高维证明需要 Vitali 型覆盖与 Lebesgue 微分工具，本书略去证明。area 与
coarea 公式还要同时处理 Jacobian 和像点重数，属于几何测度论的进一步内容。这里把
Rademacher 定理作为明确标出的外部背景记录，而不把它纳入本章已经证明的可用结论；本书后文的
任何证明也不以它为前提。

曲线长度可以直接由折线逼近定义。这个定义不依赖积分，并使可求长性成为一个纯粹的几何性质。

\begin{definition}[可求长曲线]\label{def:2-4-3-3}
设 $\gamma:[a,b]\to\R^n$ 连续。对分割
$P=\{a=t_0<t_1<\cdots<t_N=b\}$，令
\[
 L(\gamma,P)=\sum_{j=1}^N|\gamma(t_j)-\gamma(t_{j-1})|,
\]
并定义
\[
 \ell(\gamma)=\sup_P L(\gamma,P).
 \tag{2.4.23}\label{eq:curve-length-partitions}
\]
若 $\ell(\gamma)<\infty$，称 $\gamma$ 可求长。
\end{definition}

\begin{proposition}[弧长参数化]\label{proposition:2-4-3-arclength-param}
若 $\gamma:[a,b]\to\R^n$ 可求长，且 $L=\ell(\gamma)$，则存在
$1$-Lipschitz 映射 $\widetilde\gamma:[0,L]\to\R^n$，使
\[
 \widetilde\gamma([0,L])=\gamma([a,b]).
\]
原曲线停驻的参数区间不会在新参数中制造歧义。
\end{proposition}

\begin{proof}
记 $\ell(\gamma|_{[u,v]})$ 为限制曲线的长度。若 $u\le v\le w$，拼接两侧分割给出
\[
 \ell(\gamma|_{[u,w]})
 \ge\ell(\gamma|_{[u,v]})+\ell(\gamma|_{[v,w]}).
\]
反过来，在 $[u,w]$ 的任意分割中加入 $v$；三角不等式说明加点不减折线和，然后把两侧分别用其
长度控制，故反向不等式也成立。于是长度对相邻子区间可加。

令
\[
 s(t)=\ell(\gamma|_{[a,t]}).
\]
$s$ 单调不减，而且相邻区间可加性给
\[
 s(u)-s(t)=\ell(\gamma|_{[t,u]})\qquad(a\le t\le u\le b).
\]
下面证明它连续。若它在 $t$ 处不右连续，设
$\eta=\lim_{u\downarrow t}\ell(\gamma|_{[t,u]})>0$。递归选择
$u_1>u_2>\cdots>t$。给定 $u_j$，取 $[t,u_j]$ 的分割
$t=t_0<t_1<\cdots<t_N=u_j$，使折线和大于 $5\eta/6$；再利用 $\gamma$ 在 $t$ 连续，选择
$u_{j+1}\in(t,\min\{t_1,t+(u_j-t)/2\})$，使
$|\gamma(u_{j+1})-\gamma(t)|<\eta/6$。删除分割中的 $t$、加入 $u_{j+1}$ 后，三角不等式给
\[
\begin{aligned}
 \ell(\gamma|_{[u_{j+1},u_j]})
 &\ge |\gamma(t_1)-\gamma(u_{j+1})|
      +\sum_{i=2}^N|\gamma(t_i)-\gamma(t_{i-1})|\\
 &\ge \sum_{i=1}^N|\gamma(t_i)-\gamma(t_{i-1})|
      -|\gamma(u_{j+1})-\gamma(t)|
 >\frac{2\eta}{3}.
\end{aligned}
\]
此外 $u_j\downarrow t$，而这些相邻区间互不重叠。反复使用长度可加性，对每个 $N$ 有
\[
 L\ge \ell(\gamma|_{[u_{N+1},u_1]})
 =\sum_{j=1}^N\ell(\gamma|_{[u_{j+1},u_j]})
 >\frac{2N\eta}{3},
\]
这与 $L<\infty$ 矛盾。故 $s$ 右连续；反向参数化同样
证明左连续的细节如下。令
\[
 \gamma^{\leftarrow}(r)=\gamma(a+b-r),\qquad a\le r\le b.
\]
它仍连续且长度为 $L$。把已经证明的右连续结论应用于其长度函数
$s^{\leftarrow}(r)=\ell(\gamma^{\leftarrow}|_{[a,r]})$。若 $u\uparrow t$，则
$a+b-u\downarrow a+b-t$，并且
\[
 s(t)-s(u)
 =\ell(\gamma|_{[u,t]})
 =s^{\leftarrow}(a+b-u)-s^{\leftarrow}(a+b-t)\longrightarrow0.
\]
故 $s$ 左连续。

因此 $s$ 连续，并把 $[a,b]$ 映到 $[0,L]$。若 $s(t_1)=s(t_2)$ 且 $t_1<t_2$，长度可加性给
$\ell(\gamma|_{[t_1,t_2]})=0$，从而
$|\gamma(t_2)-\gamma(t_1)|=0$。所以可对 $u=s(t)$ 定义
\[
 \widetilde\gamma(u)=\gamma(t),
\]
而不依赖 $t$ 的选择。若 $u=s(t)<v=s(r)$，可取 $t<r$，并有
\[
 |\widetilde\gamma(v)-\widetilde\gamma(u)|
 =|\gamma(r)-\gamma(t)|
 \le\ell(\gamma|_{[t,r]})=v-u.
\]
故 $\widetilde\gamma$ 为 $1$-Lipschitz；$s$ 的满射性还说明两条曲线的轨迹相同。
\end{proof}

\begin{example}[折线路径的长度]\label{ex:2-4-3-polygonal}
设 $a=t_0<\cdots<t_N=b$，$\gamma$ 在每段 $[t_{j-1},t_j]$ 上线性，顶点为
$p_j=\gamma(t_j)$。含全部 $t_j$ 的分割给出
\[
 \ell(\gamma)\ge\sum_{j=1}^N|p_j-p_{j-1}|.
\]
对任意分割加入全部 $t_j$。每条线性段上按参数顺序排列的弦同向共线，其长度和等于该段两端距离；
故加细后的折线和不超过右端总和。于是
\[
 \ell(\gamma)=\sum_{j=1}^N|p_j-p_{j-1}|.
\]
弧长参数化就是以单位速度依次走过这些线段；零长度线段被自动略去。
\end{example}

\begin{proposition}[绝对连续曲线的长度]\label{proposition:2-4-3-4}
若 $\gamma:[a,b]\to\R^n$ 绝对连续，则
\[
 \ell(\gamma)=\int_a^b|\gamma'(t)|\,\dd t.
 \tag{2.4.24}\label{eq:future-length-integral}
\]
\end{proposition}

绝对连续性及 Lebesgue 微积分基本定理将在第~4~章建立，公式
\eqref{eq:future-length-integral} 的证明见\cref{proposition:4-2-3-curve-length}。上面的弧长
参数化则只使用分割上确界。

\begin{figure}[htbp]
  \centering
  \begin{tikzpicture}[node distance=7mm and 9mm]
    \node[book strong node] (lip) {Lipschitz\quad $r\mapsto Lr$};
    \node[book node,below=of lip] (lipout) {Hausdorff 维数不增\\同维零集保持};
    \draw[book accent arrow] (lip)--(lipout);

    \node[book warm node,right=of lip] (cantor) {Cantor\quad $r\mapsto r^\alpha$, $\alpha<1$};
    \node[book node,below=of cantor] (cantorout) {$C$ 零测\quad $F(C)=[0,1]$};
    \draw[book arrow] (cantor)--(cantorout);

    \node[book warning node,right=of cantor] (peano) {Peano\quad 四进制层级};
    \node[book node,below=of peano] (peanoout) {$[0,1]$ 的像\quad $[0,1]^2$};
    \draw[book arrow] (peano)--(peanoout);

    \draw[book dashed arrow] (cantor.west) to[bend right=20]
      node[above,book label] {线性尺度失效} (lip.east);
    \draw[book dashed arrow] (peano.west) to[bend right=20]
      node[above,book label] {面积被填满} (cantor.east);
  \end{tikzpicture}
  \caption{三种尺度行为。Hausdorff 覆盖估计精确表达了 Lipschitz 情形的维数控制；另外两列说明
  单独的连续性为何不足。}
  \label{fig:2-4-3-scale-comparison}
\end{figure}

局部 $C^1$ 映射在每个紧盒上导数有界，均值估计与引理
\ref{lemma:2-4-3-compact-lipschitz} 因而给出局部 Lipschitz 控制。特别地，$C^1$ 微分同胚会在
局部把零测异常集送到零测集。Brownian 路径一类远不满足此尺度控制的对象，则需要用 Hausdorff
维数描述其几何大小，而不能依赖可求长参数化。

\begin{exercise}[Lipschitz 图是零集]\label{exercise:2-4-3-1}
补全例 \ref{ex:2-4-3-2} 中底域分割的整数取整细节，并证明局部 Lipschitz 图也在每个有界柱体中
是 $n$ 维零集。
\end{exercise}

\begin{exercise}[Hausdorff 畸变]\label{exercise:2-4-3-2}
从 \eqref{eq:hausdorff-delta} 出发证明内容版本
$\mathcal H^s_\infty(f(A))\le L^s\mathcal H^s_\infty(A)$，并解释为什么证明不要求 $A$ 可测。
\end{exercise}

\begin{exercise}[Lipschitz 曲线的一维大小]\label{exercise:2-4-3-3}
若 $\gamma:[a,b]\to\R^m$ 为 $L$-Lipschitz，用等长小区间覆盖证明
\[
 \mathcal H^1(\gamma([a,b]))\le L(b-a)
\]
（采用定义 \ref{def:2-4-3-2} 的规范化）。
\end{exercise}

\begin{exercise}[空间填充与 \Holder{} 指数]\label{exercise:2-4-3-4}
设 $P:[0,1]\to[0,1]^2$ 满射且满足
$|P(s)-P(t)|\le C|s-t|^\alpha$。把参数区间等分成 $N$ 段并用像集的方形覆盖证明
$\alpha\le1/2$。
\end{exercise}

至此，集合层面的体积、线性缩放和零集控制已经建立。第~3 章将从简单函数开始构造积分，并把推前
测度、局部 Jacobian 估计与零测集保持组合成非线性换元公式。
